\section{Discussion}
\label{sec:discussion}

\subsection{Key Findings and Their Interpretation}

\textbf{Finding 1: The capacity-quality trade-off operates at 14M/31M parameter scale.}

The $\tau^*(14\text{M}) < \tau^*(31\text{M})$ result is consistent with the capacity-quality trade-off hypothesis: smaller models benefit from cleaner, lower-diversity data because their representational capacity is insufficient to extract useful signal from high-entropy distributions. Larger models leverage the distributional variety preserved by looser filtering.

\textbf{Finding 2: A minimum model scale exists for the interaction in our experiments.}

The failure of 7M/16M proxy models to exhibit any signal suggests that scale-dependent curation effects have a minimum capacity threshold. Both models must exceed this threshold, and the scale ratio must be large enough that the two models fall on opposite sides of the optimal $\tau$ curve. Whether this is a universal threshold or experiment-specific requires further testing.

\textbf{Finding 3: MMLU is unusable for sub-100M pre-training ablations.}

MMLU 4-shot accuracy equals the random baseline (0.25) for all sub-100M models at 200--500 training steps. HellaSwag 0-shot is the appropriate primary metric for pre-training ablations at this scale.

\subsection{Limitations}

\textbf{L1: PoC scale is far below target scale (14M/31M vs.\ 70M/160M).}
We confirm existence at PoC scale. The direction should hold at larger scale based on the capacity-quality mechanism, but magnitude and exact optimal $\tau$ values may differ. The full pipeline is implemented and validated (23/23 pytest tests), requiring only compute to run at 70M/160M $\times$ 50B tokens.

\textbf{L2: Effect size is small ($\Delta$\texttt{acc\_norm} $\approx 0.003$).}
At 500 training steps far from convergence, small effects are expected. Pre-training data quality effects amplify with training duration~\cite{Zhou2024ProX,Nguyen2025REWIRE}.

\textbf{L3: Deduplication $\times$ scale interaction not analyzed (P3 inconclusive).}
The data exists in \texttt{results.csv}; a post-hoc analysis from existing data requires approximately one hour. Reserved for future work.

\textbf{L4: Learning curves unavailable (P4 inconclusive).}
Disk space constraints required checkpoint deletion after each evaluation. Future runs should maintain $\geq$5 intermediate checkpoints.

\textbf{L5: Single architecture and corpus.}
All experiments use Pythia architecture and FineWeb. Standard caveat in pre-training ablation literature~\cite{Zhou2024ProX,Wettig2025Organize}.

\subsection{Broader Impact}

This work improves efficiency of language model pre-training by enabling scale-specific data curation, potentially reducing compute waste from suboptimal data recipes. No direct negative impacts are identified.
